\begin{python0}
from solutions import *; clear()
\end{python0}
\sectionthree{Why is it important? Polymorphism}

Whenever you have an object with methods with a similar name and you can't decide which one method to call until during runtime, then you will need polymorphism.

Suppose you have an array of pointers to the \verb!BasicSpaceship! class.
Also, \verb!SpecialSpaceship! is a subclass of \verb!BasicSpaceship!. In both classes you have an \verb!update()! method. The \verb!update()! method is different for the two classes: the method basically moves the objects to a new position on the screen.

\begin{console}
BasicSpaceShip * s = new BasicSpaceShip*[100];
int M = randint(); // random int from 0 to 99
for (int i = 0; i < M; i++)
{   
    s[i] = new SpecialSpaceShip();
}

for (int i = M; i < 100; i++)
{   
    s[i] = new BasicSpaceShip();
}

while (1)
{   
    for (int i = 0; i < 100; i++)
    {   
        s[i]-> update();
    }
}
\end{console}

Note that in the while-loop:

\begin{console}
while (1)
{   
    for (int i = 0; i < 100; i++)
    {
        s[i]-> update();
    }
}
\end{console}

You do not know what type of spaceship is \verb!s[3]! and therefore you do not know which update method was invoked. The update method is selected during runtime.

Why is this important? Because if one day you want to add another type of spaceship:

\begin{console}
BasicSpaceShip * s = new BasicSpaceShip*[100];
int M = randint(); // rand int from 0 to 99
for (int i = 0; i < M/2; i++)
{   
    s[i] = new SpecialSpaceShip();
}
for (int i = M/2; i < M; i++)
{   
    s[i] = new BasicSpaceShip();
}
for (int i = M; i < 100; i++)
{   
    s[i] = new VeryBasicSpaceShip();
}
\end{console}

This means that the while loop does not need to change:

\begin{console}
while (1)
{   
    for (int i = 0; i < M; i++)
    {   
        s[i]-> update();
    }
}
\end{console}

In fact \verb!s[3]! can even point to something different within the loop
and everything still works:

\begin{console}
while (1)
{   
    for (int i = 0; i < M; i++)
    {   
        s[i]-> update();
    }
    delete s[3];
    switch (rand() % 3)
    {   
        case 0: s[3] = new SpecialSpaceShip(); break;
        case 1: s[3] = new BasicSpaceShip(); break;
        case 2: s[3] = new VeryBasicSpaceShip(); break;
    }
}
\end{console}

Or for instance in your simulation, if \verb!s[3]! is initially pointing to a \verb!SuperShip object! but when hit, it points to a \verb!VeryBasicShip! object.

As long as \verb!s[3]! points to a subclass of BasicSpaceShip that has an update() method, the while loop will work.

In general:

\begin{itemize}
\item
  You analyze a system and look for objects and classify them into
  classes.
\item
  You then study what the objects do -- as general as possible within
  the class.
\item
  Design the most general parent classes. Write down relevant methods in
  parent classes and write specific methods in the children classes.
\end{itemize}

In your main program:

\begin{itemize}
\item
  In the initialization, you might have pointers (or references) to
  refer to specific children
\item
  After the initialization code, refer as much as possible to the
  methods in the parents
\end{itemize}

Going back to the simple game. Suppose game objects (spaceships) move about on their own, and dies when they collide. Thinking as generally as possible, the while loop might look something like this:

\begin{console}
while (1)
{   
    for (int i = 0; i < M; i++) ship[i]-> move();
    for (int i = 0; i < M; i++)
    {
        for (int j = i+1; j < M; j++)
        if (ship[i]-> collideswith(ship[j]))
        {   
            ship[i]-> dies();
            ship[j]-> dies();
        }
    }
    for (int i = 0; i < M; i++) ship[i]-> draw();
} 
\end{console}

This tells you that the parent class must have the following methods which the specific child classes must overwrite:

\begin{itemize}
\item
  \verb!move!
\item
  \verb!collideswith!
\item
  \verb!dies!
\item
  \verb!draw!
\end{itemize}

In the initialization, as long as the pointers \texttt{ship[i]} point to a class that subclasses the parent class, the program would work. In particular

\begin{itemize}
\item
  You can switch out subclasses
\item
  You can add new subclasses (as long as they subclass the parent class and has the above four methods)
\end{itemize}

This is a simple example where there' s only one main
parent class. In general there are (of course) many parent classes. For more complicated examples you would have to do into OO analysis/design and design patterns. Take CISS438.

Note from the above example, the methods of the parent class (example: draw) might do nothing:

\begin{console}
class BasicSpaceShip
{
public:
     void draw() {} // empty body
};

class SpecialSpaceShip
{
public:
     void draw() { /* actual drawing code */ }
};
\end{console}

This might be the case where every subclass of \verb!BasicSpaceShip! has their own \verb!draw()! method, which means that the \verb!draw()! of \verb!BasicSpaceShip! class is never used anyway.

The only purpose of the draw method in the parent is to allow polymorphism to select the child' s draw method. There' s a way to tell C++ that a method in a parent is meant to be empty and to be fulfilled by a child -- see the next section on pure virtual method and abstract base class. Abstract base class are abstract in the sense that you cannot use them to instantiate objects. Their purpose is only to act as a parent for subclasses so that you can perform polymorphism.

This style of programming is called polymorphism. One would say the code is polymorphic. Or, if the pointer points to the child but was declared as a pointer to a parent class:

\tab[5em]{\verb!P * p;!}\\

one would say that is polymorphic, or the object pointed to is polymorphic, i.e., sometimes \texttt{p} points to a class \texttt{C0} object and sometimes it points to a class \texttt{C1} object.

Some people also say the class in this case is a polymorphic type.

Don' t forget that I mentioned earlier that you can also use reference for polymorphic programming:

\tab[5em]{\verb!P & ref = c0;!}\\

In this case you would say \verb!ref! is polymorphic.

